%% kapa: exam=1 round=36 year=2025 period=2 form=A subject=accounting subject_ko=회계학 seq=14 number=54
%% source: Q-Net 감정평가사 기출문제 (cst00309) artlSeq=5231525 file=raw/1008/5231525/2025년 제36회 감정평가사 1차시험 2교시 문제지 원본.pdf
%% conversion: kordoc 4.24.51 -> tools/gen1.py (auto); review=proofread-sonnet; flags=none
\begin{question}{54}
(주)감평은 20×1년 초 (주)한국이 3년 만기로 발행한 사채(발행일 20×1년 초, 액면금액 ￦100,000, 표시이자율 연 10 \%, 매년 말 이자 지급)를 발행일의 공정가치인 ￦105,151에 취득하였다. 동 사채의 취득목적은 원리금을 수취하면서 매도할 목적으로 (주)감평은 동 사채를 기타포괄손익-공정가치측정 금융자산으로 분류하였다. 취득당시 금융자산에 적용된 유효이자율은 연 8 \%이다. 동 사채의 20×1년 말과 20×2년 말 공정가치는 각각 ￦100,000과 ￦95,000이다. (주)감평이 20×3년 초 동 사채를 공정가치인 ￦95,000에 매각하였다면, 동 금융자산 처분과 관련한 회계처리가 (주)감평의 20×3년 당기순이익에 미치는 영향은? (단, 화폐금액은 소수점 첫째자리에서 반올림하며, 단수차이로 인한 오차가 있으면 가장 근사치를 선택한다.)
\begin{choices}
  \item ￦1,437 감소
  \item ￦1,437 증가
  \item ￦3,563 감소
  \item ￦6,848 감소
  \item ￦6,848 증가
\end{choices}
\end{question}
\answer{4}
